% femtolisp 笔记
% license Usr
% type Note

\subsection{类型}
\begin{itemize}
\item \href{https://code.google.com/archive/p/femtolisp/wikis/Manual.wiki}{语法手册（非官方）}
\item \verb`(typeof 1)` 是 \verb`fixnum`
\item \verb`(typeof 'abc)` 是 \verb`symbol`， 又如 \verb`'+`
\item \verb`'abc` 是 \verb`(quote abc)` 的语法糖
\item \verb`(typeof 1.23e4)` 是 \verb`double`
\item \verb`(typeof "abc")` 是 \verb`array byte`
\item \verb`(typeof #t)` 是 \verb`boolean`， 还有 \verb`#f`
\item \verb`(typeof '())` 是 \verb`null`
\item \verb`(typeof (cons 1 2))` 是 \verb`pair`
\item \verb`(car 变量)` 和 \verb`(cdr 变量)` 分别取第 1 和 2 个元素
\item \verb`(cadr 变量)` 等效于 \verb`(car (cdr 变量))`，以此类推。
\item \verb`(list 1 2 3)` 等效为 \verb`(cons 1 (cons 2 (cons 3 '())))`， 本质上是单链表。 类型仍然是 \verb`pair`。
\item \verb`(typeof (vector 1 2 3))` 是 \verb`vector`
\item \verb`(aref (vector 10 11 12) 0)` 获取第 0 个元素 \verb`10`

\end{itemize}

\subsection{计算器}
\begin{itemize}
\item \verb`(+ 1 2 3)` 加法， 另有 \verb`-, *, /`
\item \verb`(define a 1)` 定义变量
\item \verb`(print 变量)` 打印 \verb`princ` 也可以。
\item \verb`(define a 123)` 然后 \verb|`(1 ,a)| 会得到 \verb`(1 123)`。
\item \verb`(define xs '(1 2 3))` 然后 \verb``
\item \verb`(cons 1 2)` 显示为 \verb`(1 . 2)`， \verb`(cons 1 (cons 2 (cons 3 4)))` 显示为 \verb`(1 2 3 . 4)`
\item \verb`(list 1 2 3 4)` 显示为 \verb`(1 2 3 4)`
\item \verb`(eq 1 1)` 检查是否相等， 返回 \verb`#t`
\item \verb`(= 3 3)` 等效吗？
\end{itemize}

\subsection{控制流}
\begin{itemize}
\item \verb`(if (> 2 1) 'big 'small)` 返回 \verb`'big`
\item \verb`(cond  ((= x 1) 'one)  ((= x 2) 'two) ((= x 3) 'three)  (else 'big))`
\end{itemize}

\subsection{函数}
\begin{itemize}
\item \verb`(define (square n) (* n n))` 定义函数 \verb`(square 3)` 调用
\item \verb`(null? '())` 返回 \verb`#t`
\end{itemize}
